\documentclass{article}
\usepackage{amsmath, amssymb}
\usepackage{hyperref}
\usepackage{enumitem}

\title{Graph Coloring with $\Delta + 1$ Colors: Algorithms, Proofs, and Complexity}
\author{}
\date{}

\begin{document}

\maketitle

\section{Introduction}
An overview of a graph coloring algorithm for a graph with a maximum degree $\Delta$ can color using at most $\Delta + 1$ colors, assuming the EREW PRAM model.

\section{Algorithms}
\subsection{Coloring Rooted Trees}
The 6-Color-Rooted-Tree algorithm reduces the number of colors in a rooted tree iteratively:

\begin{enumerate}
    \item \textbf{Initialization}:
    \begin{itemize}
        \item Assign each node in the tree a unique initial color equal to their IDs. It is essential for the IDs to be distinct numbers from $0$ to $N - 1$, where $N$ is the size of the graph.
        \item Each node knows the size of the graph, its ID, and $\Delta$.
    \end{itemize}
    \item \textbf{Color Refinement}:
    \begin{itemize}
        \item Repeat until each node has at most $6$ distinct colors:
        \begin{itemize}
            \item Each node communicates with its parent.
            \item The node compares its color with its parent’s color bit by bit to find the smallest bit position where they differ.
            \item It creates a new color by concatenating this bit position and the value of the differing bit.
        \end{itemize}
    \end{itemize}
    \item \textbf{Termination}:
    \begin{itemize}
        \item When the number of colors is reduced to $6$, stop the iterative process.
        \item The tree is now 6-colored.
    \end{itemize}
\end{enumerate}

The above can be refined further into a 3-Color-Rooted-Tree algorithm:
\begin{enumerate}
    \item \textbf{Reduce to Three Colors}:
    \begin{itemize}
        \item Take the 6-colored tree and perform three sequential steps to reduce the number of colors to 3.
        \item In each step, recolor each node by checking the colors of its parent and children and choosing a new color that avoids conflict.
    \end{itemize}
\end{enumerate}

\subsection{Coloring Constant-Degree Graphs}
The Color-Constant-Degree-Graph algorithm extends tree-coloring to general graphs with degree $\Delta$:

\begin{enumerate}
    \item \textbf{Partition into Pseudoforests}:
    \begin{itemize}
        \item Identify a maximal pseudoforest within the graph.
        \item Remove the edges of the pseudoforest from the graph.
        \item Repeat this process until all edges are removed, partitioning the graph into $\Delta$ or fewer pseudoforests.
    \end{itemize}
    \item \textbf{Color the Pseudoforests}:
    \begin{itemize}
        \item Use the 3-Color-Rooted-Tree algorithm to color each pseudoforest in parallel.
    \end{itemize}
    \item \textbf{Combine Colors}:
    \begin{itemize}
        \item Reintroduce the edges of each pseudoforest one at a time.
        \item Recolor nodes to ensure consistency with adjacent nodes.
    \end{itemize}
\end{enumerate}

\section{Theorems}
\textbf{Theorem 1:} 6-Color-Rooted-Tree Algorithm Correctness

The 6-Color-Rooted-Tree algorithm produces a valid 6-coloring of a tree in $O(\log^* n)$ time on a CREW PRAM using a linear number of processors.

\textbf{Proof:}
\begin{itemize}
    \item If node $i$ after one step of reassigning colors had the same color as its parent, then that would mean they differed on the same $i$-th bit previously.
    \item Since the last bit of the new color is the value of the differing bit, they must still be different in at least one bit.
    \item The time complexity is $O(\log^* n)$.
\end{itemize}

\textbf{Theorem 2:} 3-Color-Rooted-Tree Algorithm Correctness

\textbf{Proof:}
\begin{itemize}
    \item The number of colors is reduced to 3 by successive stages.
    \item Each stage shifts colors downward, eliminating one color.
    \item Each stage takes constant time.
\end{itemize}

\textbf{Theorem 3:} Coloring Constant-Degree Graphs

The Color-Constant-Degree-Graph algorithm colors a constant-degree graph with $\Delta + 1$ colors in $O((\log \Delta)(\Delta^2 + \log^* n))$ time on an EREW PRAM using $O(n)$ processors.

\textbf{Proof:}
\begin{itemize}
    \item The graph is partitioned into $\Delta$ or fewer pseudoforests in $O(\Delta)$ time.
    \item Each pseudoforest is 3-colored in $O(\log^* n)$ time.
    \item Combining colors requires $O(\Delta^2)$ iterations.
    \item Total complexity: $O((\log \Delta)(\Delta^2 + \log^* n))$.
\end{itemize}

\section{Complexity Analysis}
\begin{itemize}
    \item Rooted Tree Coloring: $O(\log^* n)$ time using $O(n)$ processors.
    \item Constant-Degree Graphs: $O((\log \Delta)(\Delta^2 + \log^* n))$ time using $O(n)$ processors.
\end{itemize}

\section{References}
\begin{itemize}
    \item Goldberg, Plotkin, Shannon - Parallel symmetry-breaking in sparse graphs.
\end{itemize}

\end{document}
